\chapter{The run engine: one worker, nine signals}\label{ch:engine}

\texttt{RunEngine} is a \texttt{QThread} subclass in
\texttt{src/\allowbreak{}ate/\allowbreak{}engine/\allowbreak{}run\_\allowbreak{}engine.py}. Its \texttt{run()} method is the worker: it inserts the
\texttt{test\_run} row, walks the expanded step tree, and closes the run. Everything the screen knows
about the run arrives through the nine signals the class declares.

\section{The worker's lifetime, and why it owns itself}

A \texttt{QThread} destroyed while its thread is still running makes Qt abort the process
(\texttt{QThread: Destroyed while thread is still running}: SIGABRT on Linux, 0xC0000409 on Windows,
with nothing printed). The danger is real because the session drops its engine when it is finished,
and that drop is queued while the worker is still returning from \texttt{run()}. Whichever side lets go
last then destroys a running thread.

So a started engine keeps itself alive. \texttt{RunEngine.start} adds the engine to
\texttt{\_live\_engines} and connects \texttt{finished} to \texttt{\_\allowbreak{}forget\_\allowbreak{}ended\_\allowbreak{}worker}, which
waits for the thread and then removes it from the set. \texttt{stop\_\allowbreak{}live\_\allowbreak{}engines(timeout\_\allowbreak{}ms=5000)}
aborts every worker still in that set and waits for it, and it is registered with \texttt{atexit}, so a
worker running when the interpreter exits is not destroyed running. If it returns a non-empty list, the
process is about to be killed by Qt, and the test suite turns that into an error message that explains
itself (chapter~\ref{ch:failed-run}).

Two rules follow, and both are easy to break by accident:

\begin{itemize}
  \item Do not \texttt{parent} a \texttt{RunEngine} to a widget in the hope that Qt will clean it up.
    A parented engine is still in \texttt{\_live\_engines}, and the widget can be deleted underneath a
    live worker.
  \item Do not call \texttt{engine.wait()} on the main thread without first calling
    \texttt{request\_abort()}. \texttt{stop\_and\_wait} does both: it sets the abort flag, then waits
    \texttt{timeout\_ms}. \texttt{wait()} releases the GIL while it waits, so the worker keeps working.
\end{itemize}

\section{What the worker does, in order}

\texttt{run()} is a guard. It calls \texttt{\_execute()} inside \texttt{try/except}, and an exception
that escapes is turned into \texttt{\_\allowbreak{}end\_\allowbreak{}unexpectedly(trace)}: the run is recorded as
\texttt{INTERRUPTED}/\texttt{FAIL}, the message is printed to stderr, the application log gets it
through the \texttt{on\_error} callable, every open Device is closed, and \texttt{FINISHED} is emitted
anyway. Each of those steps is attempted separately, because the exception may have come from the store
itself. Post-test Actions do not run after an unexpected end, as they would not after a crash.

This guard matters more since packet 32a: under PyQt6 an exception leaving \texttt{QThread.run} aborted
the process; under PySide6 it prints \texttt{Error calling Python override of QThread::run()} and the
thread simply ends. Without the guard, the run row would stay open and the screen would wait for a
\texttt{FINISHED} that never arrives. \texttt{tests/\allowbreak{}test\_\allowbreak{}unexpected\_\allowbreak{}run\_\allowbreak{}end.py} proves the run is
recorded \texttt{INTERRUPTED} and \texttt{FINISHED} is still announced.

\texttt{\_execute()} then does the work in this order:

\begin{enumerate}
  \item build a \texttt{LuaHost} for the schedule, which loads and validates it, expands
    Sub-schedules, and registers the Device drivers;
  \item insert the \texttt{test\_run} row and emit \texttt{run\_uuid\_reported};
  \item run every Pre-test Action, in Pre-test order, recording an \texttt{action} row for each; a
    Pre-test Action that raises ends the run \texttt{PRETEST\_FAILED} and records every step NOTRUN;
  \item walk the step list: header rows record their place in the tree and never run; disabled steps
    record NOTRUN with no Attempt; the abort flag is checked at every step boundary;
  \item for each enabled step, run its action closure once per Attempt, recording one \texttt{attempt}
    row each. A step's result is its last Attempt's; a step whose action raised records an
    \texttt{ERROR} Attempt; a step with no report records \texttt{NOREPORT};
  \item pause on a Stop-on-Fail failure until the Operator answers (section below);
  \item run the Post-test Actions, which run even after an abort (spec story 84), each ending on a second
    abort;
  \item record the Media Release chain and every file the run used, close open Devices, and write the
    run's result and End Reason.
\end{enumerate}

\section{The nine signals}

Every one of these is emitted on the worker and delivered to the main thread, where the receiver's own
event loop runs the slot. \texttt{tests/\allowbreak{}test\_\allowbreak{}signal\_\allowbreak{}threads\_\allowbreak32a.py} drives a real run with a
Stop-on-Fail pause, a POLL and a blocking input, and records the thread each slot ran on: they all ran
on the main thread. The same test connects one signal with \texttt{DirectConnection} as a control, and
sees the worker's thread --- and planting \texttt{DirectConnection} in the main test makes it fail. A
new signal has to be added to that test.

\begingroup\small
\begin{longtable}{p{3.3cm} p{3.6cm} p{5.2cm}}
\toprule
\textbf{Signal} & \textbf{Arguments} & \textbf{When it is emitted} \\
\midrule
\endhead
\texttt{run\_uuid\_reported} & the run key & as soon as the \texttt{test\_run} row exists; also emitted
  by an unexpected end, so the report can still find the row \\
\texttt{state\_changed} & \texttt{RunStatus} & \texttt{IDLE} $\to$ \texttt{RUNNING} $\to$ \texttt{FINISHED} \\
\texttt{input\_asked} & node seq, Attempt no, input index & a blocking \texttt{ATE.input(n)} has started
  waiting; announced from \texttt{\_on\_asked} while the worker sits in \texttt{get\_raw} \\
\texttt{step\_finished} & node seq & this step's Attempts are done, so its POLL inputs close \\
\texttt{failure\_prompt} & node seq, Attempt no & a Stop-on-Fail failure is waiting for a choice \\
\texttt{prompt\_answered} & node seq, Attempt no, choice & the pause is over: emitted when answered, and
  again when the worker moves on \\
\texttt{action\_started} & \texttt{"pre"}/\texttt{"post"}, completed, total & a Pre- or Post-test Action
  is about to run, so the status bar can say \texttt{Pre-test n/total} \\
\texttt{step\_started} & node seq & an Attempt is about to run, again on each Retry \\
\texttt{log\_lines\_recorded} & the run key & as soon as a batch of log rows lands in the store, so a
  layout can paint the newest lines while the worker is still running (packet 40a) \\
\bottomrule
\end{longtable}
\endgroup

\texttt{prompt\_answered} is emitted twice on purpose (packet 25b review): the screen clears its prompt
as soon as the answer is given, and a \texttt{failure\_prompt} still queued behind it must not leave an
answered prompt painted.

Two types of argument deserve a warning. \texttt{Signal(int)} in PySide6 is a C \texttt{int}:
\texttt{emit(2**31)} with a slot connected returns normally and the slot receives -2147483648, with a
\texttt{RuntimeWarning}; \texttt{Signal(str)} turns a non-string into \texttt{''}. A slot cannot see a
bad value, so the check is on the sending side, and
\texttt{tests/\allowbreak{}test\_\allowbreak{}signal\_\allowbreak{}threads\_\allowbreak32a.py} pins every value the engine sends. Everything the engine
emits is bounded by the schedule (node sequence, Attempt number, input index) or is a string that is
always a string (\texttt{"pre"}/\texttt{"post"}, \texttt{CONTINUE}/\texttt{RETRY}/\texttt{ABORT}, a
path). Anything that could exceed \texttt{2**31 - 1} must be declared \texttt{Signal(object)}.

\section{Blocking, and how a block lets go}

Three things block the worker, and each has to give way to an abort:

\begin{itemize}
  \item \textbf{A Stop-on-Fail pause.} \texttt{\_\allowbreak{}pause\_\allowbreak{}for\_\allowbreak{}operator} blocks with no timeout until the
    Operator answers. Packet 32b found the bug that this is about: a run-level abort that arrived before
    the pause was entered used to be lost, because the pause never read the flag. It now ends on the
    abort. A real answer given at the same moment is not overwritten, and \texttt{choose} returns
    \texttt{None} once the flag is set, so the order of the two decides rather than the worker's 50 ms
    poll.
  \item \textbf{\texttt{ATE.input(n)}.} A blocking input waits in \texttt{OperatorInput.get\_\allowbreak{}raw}, which
    takes the abort flag and gives up. The ask is announced with \texttt{input\_asked} first, so a layout
    can put a widget under the step and the Operator can answer while the worker sits.
  \item \textbf{\texttt{ATE.sleep(seconds)}.} The wait is sliced, so the abort lands within about 50 ms.
\end{itemize}

A busy loop is interrupted by the Lua count hook, which calls \texttt{aborted()} every 1000
instructions and raises \texttt{aborted} once the flag is set; \texttt{pcall} and \texttt{xpcall} are
wrapped to re-raise it so an abort cannot be caught. A Device call is interrupted by
\texttt{ate/devices/abort.py} (chapter~\ref{ch:devices}).

Three things cannot be interrupted at all, and the manual says so rather than promising otherwise: a
SQLite commit already in progress, one Lua C call such as a pattern match on a huge string, and
\texttt{set\_visa\_attribute} or \texttt{resource.timeout} on the worker.

\section{What may run on the worker}

Allowed, because the abort flag is checked inside it and it never touches a widget:

\begin{itemize}
  \item a Test Step or Action closure, synchronously. \texttt{run\_engine.py}'s docstring assumes
    this: the sandbox does not remove coroutines, but spec story 83 records a busy loop in one as an
    exception to the abort bound, so \texttt{ATE.report} is attributed to a step by snapshotting the
    host's report log before the closure runs and diffing it afterwards;
  \item store writes. Each is its own commit (the Engine is autocommit), which is what makes the record
    append-only and crash-safe, and also why a run with many rows costs many commits --- see
    chapter~\ref{ch:storage};
  \item Device reads, writes, opens and closes, through the registry, which hands every call to
    \texttt{run\_abortable};
  \item the drains: \texttt{\_\allowbreak{}flush\_\allowbreak{}device\_\allowbreak{}warnings}, \texttt{\_flush\_equipment} and
    \texttt{\_record\_files} run on the worker and write what the host collected.
\end{itemize}

Not allowed:

\begin{itemize}
  \item touching a widget or importing \texttt{QtWidgets};
  \item a wait that does not check the abort flag. Every loop a worker enters must have an exit the
    abort can reach;
  \item emitting from the main thread into the engine as if it were the worker. The main thread calls
    \texttt{choose()} and \texttt{request\_abort()}; it does not step the schedule.
\end{itemize}

One detail that trips people reading log stamps: \texttt{record\_log} takes \texttt{time.monotonic()}
when the line is recorded, and the drain only maps that to wall-clock. The \texttt{at} column promises
when the worker logged the line, not when the worker got there on time. Packet 31e's stamp test was
comparing it with a deadline and was measuring the wrong thing.

\cando{trace a run from \texttt{start()} to \texttt{FINISHED}, name each signal and what blocks the
worker, and say what an abort can and cannot interrupt.}
